\documentclass[french]{article}
\usepackage[T1]{fontenc}
\usepackage[utf8]{inputenc}
\usepackage{lipsum}
\usepackage{lmodern}
\usepackage{geometry}
\usepackage{babel}
\usepackage{graphicx}
\usepackage{lastpage}
\usepackage{ragged2e}
\usepackage{float}

\geometry{
	a4paper,
	total={210mm,297mm},
	left=20mm,
	right=20mm,
	top=20mm,
	bottom=20mm,
}

\usepackage{fancyhdr}
\pagestyle{fancy}

\lhead{Champion, Loiseau\\Rochat, Schubert}
\chead{Mobile UI}
\rhead{\today}
\renewcommand{\headrulewidth}{0.4pt}
\renewcommand{\footrulewidth}{0.4pt}

\begin{document}
	
	% Page de garde
	\title{Projet Rady - Mobile UI}
	\author{Projet de groupe\\
		Champion, Loiseau, Rochat, Schubert\\
		Resp. René Rentsch\\
		HEIG-VD}
	\date{\today}
	\maketitle
	\vspace{5cm}
	\centering
	\includegraphics[scale=0.3]{../logo/icone}
	\thispagestyle{empty}
	\newpage
	
	% Table des matières
	\tableofcontents
	\newpage
	
	\justify
	
	\section{Introduction}
		 % TODO expliquer le but du document
		 
	\section{Technologies}
		% TODO expliquer quels outils sont utilisés (ionic, leaflet?, etc.)
		
	\section{Pages}
		% TODO expliquer le concept de "page"
		% TODO montrer les transitions ?
		% TODO formaliser chaque page présentes dans l'app
		
		\subsection{\texttt{splashscreen}}
			Cette page est le point d'entrée de l'application, elle permet à celle-ci de charger la page suivante en arrière-plan. Ionic2 nous offre deux manières d'afficher un splashscreen:
			\begin{enumerate}
				\item avec le plugin natif \texttt{Splashscreen}, qui utilise une ressource et est assez limité dans ses possibilités
				\item avec une page HTML/CSS qui affiche une image, beaucoup plus souple au niveau de l'affichage
			\end{enumerate}
			La solution (2) a été ici retenue pour ses plus grandes possibilités d'évolution.
		
	% TODO autres sections? importer le user-flow?
			
\end{document}
